% Section 4: Phase 3 -- Sawin multi-quadratic construction (finite parts)

Phase~3 implements the finite, explicitly checkable parts of the
\emph{simplified} construction of Section~2.1 of
\cite{Sawin2026Remarks}, which uses an unramified pro-$2$ tower over a
multi-quadratic base field rather than the original GPT pro-$3$ tower
over a cyclic cubic. The simplification is due to Sawin and yields the
explicit numerical bound evaluated in \S\ref{sec:eq22}.

\subsection{The Sawin parameters, verbatim}

The construction fixes
\begin{align*}
  T   &= \{3,\, 5,\, 7,\, 11,\, 13,\, 17\}, \\
  S   &= \{101,\, \infty\}, \\
  L_T &= \Q(\sqrt{5},\, \sqrt{13},\, \sqrt{17},\, \sqrt{21},\, \sqrt{33}),
\end{align*}
with $[L_T : \Q] = 2^5 = 32$, and sets $K = L_T(i)$ as the CM field.

These values appear verbatim in the artifact:
\begin{itemize}[leftmargin=*]
  \item \texttt{T\_SET = (3, 5, 7, 11, 13, 17)}
  \item \texttt{S\_SET\_SPLIT = (101,)} (the infinite place is implicit)
  \item \texttt{L\_T\_GENERATOR\_RADICANDS = (5, 13, 17, 21, 33)}
  \item \texttt{SAWIN\_R\_PRODUCT = 2 \(\cdot\) 3 \(\cdot\) 5 \(\cdot\) 7 \(\cdot\) 11 \(\cdot\) 13 \(\cdot\) 17 = 510510}
\end{itemize}
and the test suite pins each constant against the published source.

\subsection{Splitting of 101 in $L_T$}

A rational prime $p$ splits completely in
$\Q(\sqrt{d_1}) \cdot \Q(\sqrt{d_2}) \cdots \Q(\sqrt{d_s})$ if and only
if it splits completely in each factor $\Q(\sqrt{d_i})$ (which, for
linearly disjoint quadratic extensions, is the definition of
completely splitting in the compositum). For odd
$p \neq d_i$, splitting in $\Q(\sqrt{d_i})$ is equivalent to
$\left(\frac{d_i}{p}\right) = 1$.

\begin{verification}\label{ver:101_splits}
The artifact evaluates the Jacobi symbol
$\left(\frac{d_i}{101}\right)$ for each
$d_i \in \{5, 13, 17, 21, 33\}$ and finds all five symbols equal to
$+1$. Therefore $101$ splits completely in $L_T$.
\textit{Source: \tname{src/erdos_ant/sawin_multiquadratic.py},
function \tname{splits_completely_in_L_T}.}
Pinned in
\tname{tests/test_sawin_multiquadratic.py::test_101_splits_completely_in_each_Q_sqrt_d}.
\end{verification}

\subsection{Galois-rank counts and Golod--Shafarevich admissibility}

Let $G = \Gal(\widetilde{\Q}_T / \Q)$ be the Galois group of the maximal
pro-$2$ extension of $\Q$ unramified outside $T$. Because $T$ contains a
prime that is $3 \pmod 4$ (e.g.\ $3$, $7$, or $11$), the $2$-rank of the
maximal elementary abelian unramified quotient is $|T| - 1 = 5$,
giving
\[
  d(G_T^\infty) = |T| - 1 = 5.
\]
Adjoining $S$ adds at most $|S| - 1 = 1$ Frobenius relations
(Sawin's pro-$2$ analogue of the Hajir--Maire--Ramakrishna cut), so
\[
  r(G_T^S) \leq d(G_T^\infty) + |S| - 1 = 5 + 1 = 6.
\]
The Golod--Shafarevich admissibility condition for the tower to remain
infinite is $r \leq d^2 / 4$; here
\[
  r(G_T^S) \leq 6 \leq \frac{(|T| - 1)^2}{4} = \frac{25}{4} = 6.25,
\]
which holds strictly. By the cited Golod--Shafarevich theorem
\cite{GolodShafarevich1964}, this admissibility implies the tower is
infinite. The artifact verifies the admissibility inequality only; it
does \emph{not} itself prove or compute the infinitude --- the cited
theorem does.

\begin{verification}\label{ver:gs_admissible}
The artifact computes $d(G_T^\infty) = 5$, $r(G_T^S) \leq 6$, and the
Golod--Shafarevich threshold $25/4 = 6.25$, and confirms
$6 \leq 6.25$. Pinned in
\tname{test_d_G_infty_equals_T_minus_1_because_T_contains_3_mod_4},
\tname{test_r_bound_is_six_per_remarks_pdf}, and
\tname{test_golod_shafarevich_inequality_holds}.
\end{verification}

\subsection{Non-claims}

Phase~3 does \emph{not} construct the infinite tower itself. The
existence of the tower is a non-constructive consequence of the
Golod--Shafarevich theorem \cite{GolodShafarevich1964} together with
Sawin's pro-$2$ analogue of the Hajir--Maire cut
\cite{HajirMaire2001,Sawin2026Remarks}; no finite computation can
exhibit a closed-form description of $L_j$ for arbitrary $j$. Phase~3
checks every step of the construction that \emph{is} amenable to finite
computation, and stops at that boundary.
